% !PRINTED 244
\weberpage{244}

The case $D=0$ likewise gives no occasion here for further distinctions; for in this case the two roots of the quadratic equation are equal to one another and real, and of the three roots of the cubic equation two are equal to one another and all three are real; for an imaginary root cannot occur doubly here, since otherwise its conjugate would also occur doubly, and just as little can the single root be imaginary, since otherwise a second one would have to be present.

For the cubic equation it only remains to seek the condition under which all three roots are equal to one another.

In this case the left-hand side of the cubic equation must be a complete cube, so that
\[
a_0 x^3 + a_1 x^2 + a_2 x + a_3 = a_0 (x - \alpha)^3.
\]
Comparison of the two sides of this identical equation yields
\[
a_1 = -3a_0\alpha, \quad a_2 = 3a_0\alpha^2, \quad a_3 = -a_0\alpha^3,
\]
from which, by elimination of $\alpha$, one obtains
\weq{1}{\begin{aligned} a_1^2 - 3a_0a_2 &= 0,\\ a_1a_2 - 9a_0a_3 &= 0. \end{aligned}}
From this there results as a consequence (by multiplying the first of these equations by $a_2$, the second by $a_1$, and subtracting)
\[
a_2^2 - 3a_1a_3 = 0;
\]
and when these conditions are satisfied, there follows from them conversely
\[
f(x) = a_0\left(x + \frac{a_1}{3a_0}\right)^3,
\]
i.\,e. the equality of all three roots.

For the cubic equation
\[
x^3 + a_1 x^2 + a_2 x + a_3 = 0
\]
we can decide completely not only about the reality but also about the signs of the roots.

Namely, if we denote by $\alpha$, $\beta$, $\gamma$ the three roots, then
\weq{2}{a_1 = -(\alpha+\beta+\gamma), \quad a_2 = \alpha\beta+\alpha\gamma+\beta\gamma, \quad a_3 = -\alpha\beta\gamma.}

Now if $D<0$, so that only one root, say $\alpha$, is real, then $\beta\gamma$ is positive and $\alpha$ will be negative or positive according as $a_3$ is positive or negative. But if $D$ is positive, so that all three roots are real, then, when $\alpha$, $\beta$, $\gamma$ are positive, by (2)
%CONT>>
